\documentclass[11pt,letterpaper]{article}
\usepackage[margin=.67in,headheight=15pt,headsep=.18in,footskip=.28in]{geometry}
\usepackage[T1]{fontenc}
\usepackage[scaled=.96]{helvet}
\renewcommand{\familydefault}{\sfdefault}
\usepackage{amsmath,amssymb,array,booktabs,fancyhdr,xcolor,hyperref}
\definecolor{ink}{RGB}{34,42,48}
\definecolor{blue}{RGB}{31,88,126}
\color{ink}
\hypersetup{colorlinks=true,urlcolor=blue,linkcolor=blue}
\pdfinfoomitdate=1\pdftrailerid{}\pdfsuppressptexinfo=15
\pagestyle{fancy}\fancyhf{}
\fancyhead[L]{\small Week 61 / Triangles on a ball / Facilitator}
\fancyfoot[L]{\scriptsize Bellingham Math Circle / Week 61 / W61-F-v1}
\fancyfoot[R]{\scriptsize\thepage}
\renewcommand{\headrulewidth}{.35pt}\renewcommand{\footrulewidth}{.25pt}
\setlength{\parindent}{0pt}\setlength{\parskip}{6pt}
\setlength{\tabcolsep}{5pt}\renewcommand{\arraystretch}{1.18}
\newcommand{\topic}[1]{\par{\large\bfseries\color{blue}#1}\par}
\newcommand{\sub}[1]{\par\textbf{#1}\par}
\newcommand{\hint}[1]{\textit{Optional hint:} #1\par}
\newcommand{\degc}{^\circ}
\begin{document}
\topic{Mathematical overview: what is true, and when}
\textbf{Model.} Use an ideal round sphere of radius $R>0$; its surface area is $4\pi R^2$. A great circle is the intersection with a plane through the center. Its arcs are locally straight \emph{on the surface}, although curved in space. Two distinct nonantipodal points have one shortest surface route: the minor great-circle arc, of length $R\arccos(u\cdot v)$ for unit position vectors $u,v$. Antipodal points have infinitely many shortest routes: great-circle \emph{semicircles} of length $\pi R$. A full great circle is neither a semicircle nor a shortest endpoint-to-endpoint route. Non-equatorial latitudes are not great circles.

\textbf{Lune and triangle area.} A lune between two great-circle semicircles with common antipodal endpoints and angle $\theta$ degrees covers $\theta/360$ of the sphere. A pole/equator triangle with meridian gap $0<\theta<180$ has corners $\theta,90,90$ and fraction $\theta/720$.

For a positive-area convex triangle $T$ with three distinct, nonantipodal, linearly independent vertices, minor great-circle sides, and the \emph{whole smaller region strictly inside some open hemisphere}, its inside corners $\alpha,\beta,\gamma$ are each between $0$ and $180$ degrees. Girard's formula is
\[
  \operatorname{area}(T)=R^2(\alpha+\beta+\gamma-\pi)\quad\text{(radians)},
  \qquad f=\frac{\text{sum in degrees}-180}{720}.
\]
Thus $180<\text{sum}<540$ and $0<f<1/2$. Corners mean tangent-direction angles on the sphere; a flat projected drawing does not measure them. The formula here does not cover long sides, latitude sides, degenerate triples, or the complementary region.

\textbf{Exact reason.} Extending the three sides to full circles makes eight nonempty cells. Select the containing and opposite lune at each corner: three pairs, six lunes. $T$ and its antipodal copy $T^*$ are covered three times; the other six cells once. The antipodal map preserves area. Therefore the summed pair areas equal one whole sphere plus four extra copies of $T$. Pages 7--8 prove every part of this count, without assuming equal cells.

\textbf{Discovery and evidence.} Children can find different shortest routes between opposite points, a three-right-corner triangle, and increasing area as a meridian gap widens. Handling and measurement are experiments. Two covering examples support a conjecture; the sign-cell argument explains why the count always persists. Measurement, exact supplied records, and a fitted formula alone are not a universal proof.

\begin{tabular}{@{}p{.82in}p{5.95in}@{}}
\toprule
Entry & Prerequisites and destination\\\midrule
K--1 & Supported route/corner exploration (page 3 of this guide); no theorem packet or independent geodesic-understanding claim.\\
P1--3, 2--5 & Adult reading; fixed endpoints, surface contact, right-angle recognition. Choose and check constructions; no angle sums required.\\
P4, 3--5 & Equal rotations, degree gaps, halves and whole-sphere fractions; adult arithmetic recording is possible.\\
P5--7, 4--5 & Track overlapping pair memberships on one ball; counted area can exceed a whole. Add through 300, subtract 180, use fractions. Universal explanation by readiness.\\\bottomrule
\end{tabular}

\newpage
\topic{Preparation for eleven children and three anchored adults}
This guide uses the current roster: four K--1 children, four third graders, and two fourth graders plus one fifth grader. The parent stays with K--1; a mathematician stays with Grades 2--3; the organizer stays with Grades 4--5. Keep tables fixed. Two pairs work at each younger table; the upper trio rotates builder, checker and recorder. The launch uses the upper kit, so it needs no sixth ball.

\begin{tabular}{@{}p{1.04in}p{.56in}p{5.17in}@{}}
\toprule
Table & Kits & Actual allocation\\\midrule
K--1 & 2 & Two plain firm 15--20 cm balls, two supports, six 80 cm strings, twelve reusable point dots, four narrow paper right angles.\\
Grades 2--3 & 2 & Two suitable equator/pole models, two supports, six strings, twelve dots, four right angles.\\
Grades 4--5 & 1 & One equator/pole model, one support, three strings, six dots, three right angles. Same model serves whole-group launch and octant covering.\\\bottomrule
\end{tabular}

\textbf{Total:} five balls/models and five stable bowls or rings; fifteen strings (12 m total); thirty reusable small dots plus ten spares; eleven narrow right-angle tabs. Prepare tabs from three scrap Letter sheets. Bring eleven pencils, three washable fine markers and three wiping cloths (one of each per table), plus one setup ruler and scissors. Each triangle has \emph{three} recoverable sides: use three strings, or retain the third side as a removable verified arc. Do not erase a side while checking its corners.

\textbf{Printing.} First visit: P1--3 for seven older children = 21 single-sided Letter sheets; four plain sheets for K--1 drawings. Give one page at a time. Return menu: P4 for seven children = 7 sheets; P5--7 for the upper trio = 9 sheets. The full menu is 37 student sheets plus 4 plain sheets, across visits. Three copies of this nine-page guide = 27 adult sheets. Print at 100\%; student diagrams are records, not ball-fitting templates.

\textbf{One-time model setup.} Three older-table models need a checked equator, poles $N,S$, and perpendicular meridians for a known right-angle example and octant. A meridian is a pole-to-pole great-circle semicircle; the continuation past $S$ completes its full circle. For P4 add removable meridian references at $45,72,120$ degrees from one baseline, when needed. Use a manufactured calibrated globe or center-plane/angular reference; a protractor on a projected picture or laid across a curved patch is insufficient. Keep only the currently used references visible. P6's $80\degc$ model and P7's records are supplied on paper; do not budget exact physical fabrication of those triangles.

\textbf{Physical pretest before classroom use.} With these exact kits, test support stability; child-reachable surface contact with both endpoints fixed; retaining three arcs; locating opposite dots; showing back arcs by turning the ball; and comparing tangent directions at a \emph{known} right corner with a narrow paper tab. Verify every intended meridian gap, complete-circle marking, containing/opposite lune choice and eight-region record. Test that marks remove. A taut friction-held loop can remain on a small circle; string tension alone is not calibration. An 80 cm string exceeds the 63 cm circumference of a 20 cm ball, but knots, grip and handling still need testing.

\textbf{Status and forecast.} Fabrication, calibration, physical fit, string/corner/lune handling rehearsal and classroom piloting are \textbf{UNPERFORMED}. With suitable models in hand, allow about 60--90 minutes initially for marking, cutting and rehearsing; sourcing/calibration may take longer. Reuse may take 10--20 minutes to inspect, distribute and refresh marks; P4 references can add 15--25 minutes. These are untested forecasts, not a measured preparation time.

\newpage
\topic{First visit: launch together, then retain children's choices}
\begin{tabular}{@{}p{.65in}p{6.12in}@{}}
\toprule
Minutes & Suggested hour; finishing a packet is not the target\\\midrule
0--5 & Movement, then move to the three fixed tables.\\
5--8 & Let pairs/trio handle balls and string on supports. Notice contact, rolling and the hidden side before instructions.\\
8--12 & Brief shared legal-action demonstration below.\\
12--29 & K--1 route/corner play; older tables P1 and chosen P2 attempts.\\
29--45 & Continue P2 or begin P3; let a promising construction have time.\\
45--54 & Repeat/compare constructions. P4 only if the table already controls the model and equal slices.\\
54--60 & Share one route, corner or surprising triangle. Save a drawing and actual stopping point for the next visit.\\\bottomrule
\end{tabular}

\textbf{Launch prompt.} ``Keep these two dots fixed. Make a route that touches the ball all the way. Change the route while the dots stay put.'' Let a child move the string; show a lifted chord beside a surface route. Use P1's existing $P,Q\to$ full circle $\to$ shorter arc visual for a practice pair. Say that the ideal straight surface route follows a \emph{great circle}; trace its shorter arc on the checked model. Draw the same two labels and arc as one legal record; dashed means behind the ball. Do not solve the opposite-point question or show the octant yet.

Hold two short side directions at a known corner; compare with the narrow paper right angle. Use P2's existing $P$ example: two surface directions, their local $135\degc$ drawing, matching angle record. Children need to know what is being compared, not add 135 to anything. Adult checks the hemisphere/convexity rule while children choose points. Reading rules or steadying a support is ordinary assistance; choosing all points and routes for a child operates the investigation for them.

\textbf{K--1 concrete route.} In each pair, one child chooses two visible dots and holds an endpoint, the other makes a contacting string route; swap. Ask them to make a different route between the \emph{same} dots and show which seems shorter. They may draw one route on plain paper or ask the parent to record their choice. Then put two short routes at one shared dot and find a corner that matches the paper right angle, and one that does not. Adults can hold the ball, but children choose what changes. This is substantial object exploration; it does not claim a shortest-path proof, degrees, area fractions or K--1 coverage by P1--7. Repeat and vary it while the older tables investigate triangles. If a pair cannot keep the rule after the launch, record that operational limit rather than forcing a theorem task.

\textbf{Return gates.} P4 needs children to keep poles/equator fixed and understand equal rotations. P5 needs a successful single-pair record; P6 needs a recoverable eight-region membership row. P7 needs children to ask why examples generalize and to accept counted areas above one sphere. These are readiness checks, not grade tests. The anchored upper adult may explain fractions while children make covering choices; do not silently replace their choices with an adult tally.

\textbf{Later covering launch.} Reuse P4's non-task $40\degc$ lune $\to$ nine rotations $\to1/9$ record before area work. Before P5, reuse its non-task $80\degc$ example: corner $A\to$ circles $AB,AC\to$ containing/opposite pair $\to X$ counted once. Rotate or point to the back of the \emph{same} ball. One selected pair contributes 0 or 1, never 2, to an interior patch. An unmarked pair contributes 0 to the current record. The full three-pair count remains the children's task.

\newpage
\topic{Problems 1--3: routes, corners and a variable triangle}
\sub{Problem 1 / Grades 2--5 / chosen close, far and opposite pairs}
\textbf{Answer.} Both distinct nonopposite pairs have a unique shortest route, their minor great-circle arc. Only the opposite pair permits two different routes of the same shortest length; in fact infinitely many semicircles work. Close/far is not the uniqueness distinction. The initial $P,Q$ example is a $95\degc$ minor arc; the other $265\degc$ arc of that circle is longer.

Let children compare several contacting routes before naming the theorem. A string that lifts off gives a chord, a different problem. A route along a non-equatorial latitude may look straight in one view but is not a great-circle route. ``Meridian'' refers to one pole-to-pole semicircle; extending beyond the opposite pole gives the full circle.

\hint{On a checked model, keep the endpoints and turn the ball so you can see the circle containing them. For opposite dots, try passing through two different third points.}
\textbf{Adult reason.} Rotate the sphere so the first endpoint is the north pole. For colatitude $t$ and longitude $\phi$, surface length satisfies $ds^2=R^2(dt^2+\sin^2t\,d\phi^2)$, so any route to colatitude $t_0$ has length at least $Rt_0$. Equality follows one longitude with monotone colatitude. For $0<t_0<\pi$ the second endpoint fixes that longitude; at $t_0=\pi$ every longitude gives a semicircle. This establishes the distance and uniqueness for ordinary rectifiable shortest routes (equivalently after arc-length parametrization); calculus stays with adults. Petrunin--Zamora, Observation 2.22, pp. 24--25, gives the distance theorem; \S14B, p. 119, states the antipodal nonuniqueness.

\sub{Problem 2 / Grades 2--5 / as many right corners as possible}
\textbf{Answer.} All three are possible. Take $N=(0,0,1)$, $A=(1,0,0)$ and $B=(0,1,0)$ on a unit sphere. The sides are quarter-circles in three perpendicular center planes; the inward tangents at each vertex are perpendicular. The smaller triangle is one octant, fraction $1/8$, with sum $270\degc$. It fits strictly within the hemisphere centered toward $(1,1,1)$, even though $A,B$ lie on the chosen reference equator. ``Some open hemisphere'' does not mean the fixed northern hemisphere.

\hint{Offer the checked equator and two perpendicular meridians only after real attempts. Can their intersections become corners?}
An almost-right physical corner is experimental evidence. The perpendicular center planes give the exact reason. Keep the $135\degc$ recording example separate from the sought triangle; no angle arithmetic is needed to search.

\sub{Problem 3 / Grades 2--5 / fixed pole, moving equator points}
\textbf{Answer.} If the shorter equatorial gap is $0<\theta<180$ degrees, corners are $\theta,90,90$. Each meridian center plane contains the polar axis, so it meets the equatorial plane at a right angle. Wider shorter gap means a larger triangle. Children can choose, for example, $60$, $90$, $120$ degree gaps; exact labels are optional here. The printed example is $110\degc$ at $N$.

\hint{Keep one equator dot fixed and move the other within the same shorter half of the equator. Which two corners stay right?}
The chosen triangle is the northern half of the selected lune; P4 makes that area statement exact. At gap 0 it degenerates. At 180 the equator endpoints are antipodal, with no unique shorter third side; these are excluded endpoints, not extra examples of the triangle rule.

\newpage
\topic{Problems 4--5: one lune, then three opposite pairs}
\sub{Problem 4 / Grades 3--5 / pole--equator fractions}
\textbf{Key.} Use fraction of the \emph{entire surface}, not a flat projected disk.
\begin{center}\begin{tabular}{c|c|c|c}
Meridian gap & Equal lunes in a sphere & One lune & Northern triangle\\\hline
$45\degc$ & 8 & $1/8$ & $1/16$\\
$72\degc$ & 5 & $1/5$ & $1/10$\\
$120\degc$ & 3 & $1/3$ & $1/6$
\end{tabular}\end{center}
\textbf{Reason.} Rotating a lune about its poles preserves its area. Equal gaps partition the sphere into congruent strips; thus the $40\degc$ convention example has nine equal lunes and fraction $1/9$. Reflection across the equator preserves area and exchanges the northern and southern halves of the selected lune. The equator closes each half into a three-corner triangle, giving half the lune fraction. This proves $f=\theta/720$ for this family without using the general triangle theorem.

For an arbitrary angle, rotational symmetry and additivity give proportional area first for rational fractions of a full turn and then, by continuity, for all angles. Hence one lune of angle $\theta$ radians has area $2R^2\theta$. The same relation holds about \emph{any} antipodal pole pair, not only the marked north/south poles.

\hint{Use the non-task nine-strip example already on P4. Then ask what the equator does to one complete lune. Do not measure the projected shaded shape.}
Older children can place the marked gaps; the adult must have checked their exact reference. Children may find the half-lune relation by matching or rotating, with no written division algorithm requirement. Do not substitute the triangle formula as the first explanation of this problem.

\sub{Problem 5 / Grades 4--5 / octant covering}
Use the octant from P2 and extend \emph{all three sides to complete great circles}. Keep those circles on the same ball. At a chosen corner its two incident circles make four lunes; select the one containing the triangle and its opposite, not all four. These two lunes have disjoint interiors. A patch has membership 0 or 1 in this pair.

The P5 convention example uses $A,B,C$ on a non-task $80\degc$ triangle. Its selected circles are $AB$ and $AC$, not $BC$; patch $X$ is inside just the selected $A$ pair and records 1. The target octant uses labels $N,A,B$. Do not carry the example's corner names into the target key accidentally.

\begin{center}\begin{tabular}{c|cccccccc}
Octant region & 1&2&3&4&5&6&7&8\\\hline
Pairs at & $N,A,B$&$B$&$A$&$N$&$N$&$A$&$B$&$N,A,B$\\
Count &3&1&1&1&1&1&1&3
\end{tabular}\end{center}
Each of eight cells has fraction $1/8$ by the three perpendicular center planes. Each right-angle lune has fraction $1/4$, each opposite pair $1/2$, and the three pairs sum to $3/2$ spheres. The cell tally gives $(3+6+3)/8=3/2$ too. Region 1 is the target and 8 its opposite: each is $1/8$. To prepare the later argument, count one whole sphere first; regions 1 and 8 each have two \emph{additional} coverings.

\hint{Ask for the names of pairs covering one interior patch, then recover its total. Rotate the ball; the back picture is not another ball. Count areas of regions, not boundary crossings.}
This symmetric example explains an octant's area, but it does not establish that all eight cells are always equal, or prove the universal covering count. P6 deliberately removes that symmetry.

\newpage
\topic{Problem 6: unequal cells, the same covering counts}
\textbf{Use the supplied $80,80,80$ degree map.} Exact physical construction of this triangle is not required. Both views belong to one unchanged sphere. They look normally toward a side plane and away from it: front contains complete regions 1--4, back complete regions 5--8. The silhouette is the third full circle in the limb plane; visible arcs divide each hemisphere into four cells. Deliberate limb labels identify the same boundary points. Hidden arcs are omitted here, unlike the oblique introductory views; they must not be invented as front-face subdivisions.

The pair-letter row is the recoverable record. Let children identify each pair's two circles and selected opposite lunes on these maps; an adult can write letters that children choose. Then derive the count row from the letters. No simultaneous second tally or exact globe fabrication is needed.

\begin{center}\begin{tabular}{c|cccccccc}
Region &1&2&3&4&5&6&7&8\\\hline
Pair letters & $ABC$&$C$&$B$&$A$&$A$&$B$&$C$&$ABC$\\
Count &3&1&1&1&1&1&1&3
\end{tabular}\end{center}
Opposite region partners are $1\leftrightarrow8$, $2\leftrightarrow7$, $3\leftrightarrow6$, $4\leftrightarrow5$. At $A$ use circles $AB,AC$; at $B$, $AB,BC$; at $C$, $AC,BC$. The pair at each corner covers four whole cells, as its two defining circle sides agree. On the page, boundary samples are unsuitable count patches: select interiors.

\textbf{Area answer without assuming equal cells.} One $80\degc$ lune has fraction $80/360=2/9$; each opposite pair has $4/9$. Summed pair area is $3(4/9)=4/3$ spheres. If the target fraction is $f$, its antipodal triangle has the same $f$. The table says one whole sphere plus two extra coverings of each target triangle, so
\[
 \frac43=1+2f+2f=1+4f,\qquad f=\frac1{12}.
\]
This is an area count with overlap, not a physical union of area $4/3$. A union on one sphere cannot exceed a whole, but a sum of memberships can.

\textbf{Independent cell key.} Regions 1 and 8 have angles $80,80,80$ and fraction $1/12$ each. Each other cell is bounded by one of the other signed triples of vertex rays; its corner angles are one $80\degc$ and two supplements $100\degc$. Its sum is $280\degc$, giving fraction $100/720=5/36$ once the theorem is proved. Coordinate solid-angle checks independently agree. Total area is
\[
 2(1/12)+6(5/36)=1,
 \qquad\text{one pair}=2(1/12)+2(5/36)=4/9.
\]
These cell areas are an adult check, not needed for the student's deduction; using the theorem to get them cannot be offered as an independent proof of that theorem.

\hint{If children assume ``eight cells means eighths,'' compare the sizes of cells 1 and 2 in the same map. Ask what can be counted without knowing each area. If the one-pair meaning is lost, return to P5's $X=1$ example.}
\textbf{Discussion bridge.} The octant and this map have identical membership counts despite different cell areas. This suggests an invariant. Ask which features of the circles determine membership; page 7 supplies an exact general answer. Checking these two examples is still an experiment, not the reason all allowed triangles work.

\newpage
\topic{The universal explanation: eight cells really exist}
This is adult reasoning for P7, after concrete covering work. Keep one sphere, one triangle and three fixed side circles throughout. Let $A,B,C$ be unit position vectors of the triangle's vertices. They are linearly independent: otherwise all lie in a common center plane, and their minor-arc boundary cannot enclose the required positive-area convex triangle.

Orient a unit normal $n_A$ to the plane of $BC$ so that $n_A\cdot A>0$. Similarly orient $n_B$ toward $B$ across $CA$, and $n_C$ toward $C$ across $AB$. Each is zero on the other two vertex vectors. Call the corresponding open positive hemispheres $H_A,H_B,H_C$.

\textbf{Why exactly eight nonempty cells.} A point has a sign triple according to $n_A\cdot x,n_B\cdot x,n_C\cdot x$, away from circle boundaries. For any signs $s_A,s_B,s_C\in\{+1,-1\}$,
\[
 x=\frac{s_AA+s_BB+s_CC}{\lVert s_AA+s_BB+s_CC\rVert}
\]
exists by independence and has exactly those signs: for instance $n_A\cdot x$ has sign $s_A$, since the $B,C$ terms vanish. Thus none of the eight cells is fictitious. More generally expand any $x=aA+bB+cC$ uniquely. Its normal signs are the signs of $a,b,c$. Each sign region is the normalized positive cone generated by $s_AA,s_BB,s_CC$; it is one connected spherical triangular cell, with minor-arc edges. Hence there are eight cells, not just eight tested points. The $+++$ cell is $T$, the convex minor-arc region; the $---$ cell is its antipodal copy $T^*$. Their closures include boundaries, which have zero surface area.

\textbf{Which pair contains a cell?} At $A$, the incident circles are $AB$ and $AC$, whose normals are $n_C,n_B$. The containing lune is $H_B\cap H_C$ (with boundary), and the opposite is its antipodal image. A cell lies in that pair exactly when its $B,C$ signs agree. At $B$ compare $A,C$; at $C$ compare $A,B$. These pairs have corner widths $\alpha,\beta,\gamma$, the triangle's actual inside angles, each less than $180\degc$.

\begin{center}\begin{tabular}{c|c|ccc|c}
Region & Signs $(A,B,C)$ & Pair $A$ & Pair $B$ & Pair $C$ & Count\\\hline
1 & $+++ $ &1&1&1&3\\
2 & $++- $ &0&0&1&1\\
3 & $+-+ $ &0&1&0&1\\
4 & $+-- $ &1&0&0&1\\
5 & $-++ $ &1&0&0&1\\
6 & $-+- $ &0&1&0&1\\
7 & $--+ $ &0&0&1&1\\
8 & $--- $ &1&1&1&3
\end{tabular}\end{center}
All-equal signs give three agreements. Every other triple has exactly two equal signs and one different sign: one agreement, therefore one covering. This explains persistence for \emph{every} triangle in the stated scope, independent of symmetry and cell area. For P5 only, substitute $(A,B,C)=(N,A,B)$ to match its labels. The antipodal map $x\mapsto-x$ exchanges complementary sign triples and preserves lengths and surface area. Thus $\operatorname{area}(T^*)=\operatorname{area}(T)$ exactly; a camera's apparent sizes are irrelevant.

\newpage
\topic{Area deduction, Problem 7, and the theorem's limits}
\textbf{Count the six lune areas in two ways.} Write $S=\alpha+\beta+\gamma$ in radians. One lune of angle $\alpha$ has area $2R^2\alpha$, so its opposite pair has $4R^2\alpha$. Summing three pairs gives $4R^2S$. By the proven membership count it also gives
\[
 4\pi R^2+2\operatorname{area}(T)+2\operatorname{area}(T^*)
 =4\pi R^2+4\operatorname{area}(T).
\]
The ``4'' means two extra copies of each of two equal triangles, not four disjoint triangles. Equate, divide by 4, and obtain $\operatorname{area}(T)=R^2(S-\pi)$. Dividing by $4\pi R^2$ and converting degrees gives
\[
 f=\frac{S_{\rm degrees}-180}{720}.
\]
Positive area forces a sum above $180\degc$; the compact triangle strictly inside an open hemisphere has area strictly less than half a sphere. The formula therefore gives a sum below $540\degc$. Children may phrase the count with ``one ball plus extra copies'' instead of algebra.

\sub{Problem 7 / Grades 4--5 / rule, universal explanation and records}
\begin{center}\begin{tabular}{c|c|c|c}
Exact surface corners & Sum & Excess & Whole-sphere fraction\\\hline
$60,60,90$ & $210\degc$ & $30\degc$ & $1/24$\\
$80,80,80$ & $240\degc$ & $60\degc$ & $1/12$\\
$100,100,100$ & $300\degc$ & $120\degc$ & $1/6$
\end{tabular}\end{center}
These are supplied exact mathematical records, not requirements to construct exact $80\degc$ or $100\degc$ corners on classroom balls. Each exists under the stated rules. For the first, $A=(0,0,1)$, $B=(2\sqrt2/3,0,1/3)$, $C=(1/\sqrt6,1/\sqrt2,1/\sqrt3)$ gives those named tangent angles. For equal angle $q=80\degc$ or $100\degc$, set $d=\cos q/(1-\cos q)$ and choose three unit vectors with pairwise dot product $d$. Their Gram matrix has eigenvalues $1-d,1-d,1+2d>0$, so they are independent; $A+B+C$ has positive dot product $1+2d$ with each vertex, giving an open-hemisphere witness. Their tangent angle cosine is $d/(1+d)=\cos q$. Coordinates and solid angles are also checked separately in the source helper.

\hint{When a child proposes the rule, ask what stays true as the unequal circles move. Use the one-pair test and opposite-area reason, rather than asking them to memorize the table. If that explanation needs constant adult direction, continue P5--6 and save the universal proof for a return.}

\textbf{Counterexamples to overextension.} Three $60\degc$ corners would give zero area, so no positive-area triangle in our scope has them. Taking the octant's complementary region gives fraction $7/8$ and reflex inside corners $270\degc$; it violates our convex/open-hemisphere/less-than-$180\degc$ conditions. A long equatorial side changes the selected region and is not a minor side. A small latitude circle cannot replace the equator in P3--4: it is not cut by a center plane, and the right-angle/excess reasoning cannot be transferred to that curved-edged region.

Two distinct great circles meet at two antipodal points because their center planes meet in a line through the center; complete great circles therefore do not provide parallel lines on the sphere. Doubling $R$ doubles route lengths and quadruples absolute area, but leaves corners and whole-sphere fractions fixed. A larger physical ball's triangle can have more area without having a larger fraction.

\newpage
\topic{Return visits, teaching evidence and source provenance}
\textbf{Useful continuations.} Return first to the saved physical construction rather than begin with a formula. At the middle table, children can vary the meridian gap while keeping the two equator corners right, then explain the north/south equality. At the upper table, finish an unequal-cell pair-letter record before introducing signs. A ready child can describe a region as ``same side/different side'' of three complete circles and explain the agreement rule. Other children can keep making and checking surface triangles alongside that conversation. Another return can compare equal-angle triangles on different-radius spheres, or find which proposed angle records are impossible. No fresh student pages are needed for these adult-led extensions; use the existing pages and plain paper as desired.

\textbf{Prior use and novelty.} The atlas family GA-23 already uses the octant and variable-longitude family for parallel transport. Week 61 develops intrinsic shortest routes, surface corners, triangle area and the universal overlap count; the octant and lune relation are prior library mathematics. The Week 41 translation-glued flat torus has a different intrinsic metric: a physical ball does not model it. Returning children may know a construction without having investigated this question. Record what a child actually used rather than treating a library reference as proof of prior classroom use.

\textbf{Pedagogy actually found in sources.} Givental, Nemirovskaya and Zakharevich, \emph{Math Circle by the Bay}, printed p. viii (PDF p. 9), chooses themes with deep mathematical context; printed pp. ix--x (PDF pp. 10--11) discuss manipulatives, independent attempts, variation in pace/depth and additional challenges. Rozhkovskaya, \emph{Math Circles for Elementary School Students}, Lesson 3, ``At the lesson,'' item 1, describes attempts and explanations before an organizing table. Lesson 7, item 2, reports enthusiastic K'nex polygon construction and nonrecall of an earlier triangle experiment; Lesson 8, item 2, reports writing/coloring delays. These are particular observations, not a validation of our sphere activity or a ban on triangles.

The organizer's year-one \emph{Pattern Block Exercises -- Google Docs}, PDF p. 1, allows five minutes of free exploration before pair/trio work. Current project context reports tiling success and the failure of an abstract paper-lamp rule. \textbf{Our proposed adaptation:} allow handling first, show fixed endpoints and contact, build corners before records, externalize one membership row, and let the deeper proof occupy another visit. This is an inference from those sources. This sphere lesson and its staffing/preparation estimates remain unpiloted.

\textbf{Mathematical source.} Anton Petrunin and Sergio Zamora Barrera, \emph{What Is Differential Geometry? Curves and Surfaces}, version 7, 31 December 2025: Observation 2.22, printed/PDF pp. 24--25 (spherical distance); \S14B, p. 119 (antipodal shortest paths); Lemma A.17, p. 178 (lune covering and spherical excess); Proposition 17.10, p. 143 (spherical Gauss--Bonnet context). Primary text: \href{https://arxiv.org/pdf/2012.11814v7}{arxiv.org/pdf/2012.11814v7}. The verified primary substitute is used; no unlocated Berkeley course passage is cited. The theorem is established mathematics. The detailed sign-cell argument, guide wording, kit plan and check helper are newly authored adaptations; no source prose, figure or other asset is reproduced.

\textbf{After the first visit, record:} pages/attempts actually used by each table; whether children kept endpoints and surface contact; what they chose independently; whether corner and pair examples transferred; adult reading/holding versus adult control of choices; physical failures and fixes; actual preparation time; and the question children wanted to continue. Separate observations from guesses about their cause. Mark an unattempted continuation as untested. The source/digital checks support mathematics and print layout; they do not establish material readiness or classroom success.

\textbf{Version.} This original guide W61-F-v1 accompanies the accepted seven-page student packet W61-S-v2. It is a separate untested adult step. Portable guide source contains only this authored TeX, build/check helpers and README; downloaded references, workflow prompts, borrowed exemplars and render intermediates are excluded. No remote copy or publication permission is asserted.
\end{document}
